\documentclass[10pt, letterpaper]{article}
% Packages:
\usepackage{academicons}
\usepackage[
    ignoreheadfoot, % set margins without considering header and footer
    top=2 cm, % seperation between body and page edge from the top
    bottom=2 cm, % seperation between body and page edge from the bottom
    left=2 cm, % seperation between body and page edge from the left
    right=2 cm, % seperation between body and page edge from the right
    footskip=1.0 cm, % seperation between body and footer
    % showframe % for debugging 
]{geometry} % for adjusting page geometry
\usepackage{titlesec} % for customizing section titles
\usepackage{tabularx} % for making tables with fixed width columns
\usepackage{array} % tabularx requires this
\usepackage[dvipsnames]{xcolor} % for coloring text
\definecolor{primaryColor}{RGB}{0, 79, 144} % define primary color
\usepackage{enumitem} % for customizing lists
\usepackage{fontawesome5} % for using icons
\usepackage{amsmath} % for math
\usepackage[
    pdftitle={Hao-Hsuan Chang's CV},
    pdfauthor={Hao-Hsuan Chang},
    pdfcreator={LaTeX with RenderCV},
    colorlinks=true,
    urlcolor=primaryColor
]{hyperref} % for links, metadata and bookmarks
\usepackage[pscoord]{eso-pic} % for floating text on the page
\usepackage{calc} % for calculating lengths
\usepackage{bookmark} % for bookmarks
\usepackage{lastpage} % for getting the total number of pages
\usepackage{changepage} % for one column entries (adjustwidth environment)
\usepackage{paracol} % for two and three column entries
\usepackage{ifthen} % for conditional statements
\usepackage{needspace} % for avoiding page brake right after the section title
\usepackage{iftex} % check if engine is pdflatex, xetex or luatex

% Ensure that generate pdf is machine readable/ATS parsable:
\ifPDFTeX
    \input{glyphtounicode}
    \pdfgentounicode=1
    % \usepackage[T1]{fontenc} % this breaks sb2nov
    \usepackage[utf8]{inputenc}
    \usepackage{lmodern}
\fi



% Some settings:
\AtBeginEnvironment{adjustwidth}{\partopsep0pt} % remove space before adjustwidth environment
\pagestyle{empty} % no header or footer
\setcounter{secnumdepth}{0} % no section numbering
\setlength{\parindent}{0pt} % no indentation
\setlength{\topskip}{0pt} % no top skip
\setlength{\columnsep}{0cm} % set column seperation
\makeatletter
\let\ps@customFooterStyle\ps@plain % Copy the plain style to customFooterStyle

\makeatother


\titleformat{\section}{\needspace{4\baselineskip}\bfseries\large}{}{0pt}{}[\vspace{1pt}\titlerule]

\titlespacing{\section}{
    % left space:
    -1pt
}{
    % top space:
    0.3 cm
}{
    % bottom space:
    0.2 cm
} % section title spacing

\renewcommand\labelitemi{$\circ$} % custom bullet points
\newenvironment{highlights}{
    \begin{itemize}[
        topsep=0.10 cm,
        parsep=0.10 cm,
        partopsep=0pt,
        itemsep=0pt,
        leftmargin=0.4 cm + 10pt
    ]
}{
    \end{itemize}
} % new environment for highlights

\newenvironment{highlightsforbulletentries}{
    \begin{itemize}[
        topsep=0.10 cm,
        parsep=0.10 cm,
        partopsep=0pt,
        itemsep=0pt,
        leftmargin=10pt
    ]
}{
    \end{itemize}
} % new environment for highlights for bullet entries


\newenvironment{onecolentry}{
    \begin{adjustwidth}{
        0.2 cm + 0.00001 cm
    }{
        0.2 cm + 0.00001 cm
    }
}{
    \end{adjustwidth}
} % new environment for one column entries

\newenvironment{twocolentry}[2][]{
    \onecolentry
    \def\secondColumn{#2}
    \setcolumnwidth{\fill, 3.5 cm}
    \begin{paracol}{2}
}{
    \switchcolumn \raggedleft \secondColumn
    \end{paracol}
    \endonecolentry
} % new environment for two column entries

\newenvironment{twocolentryPub}[2][]{
    \onecolentry
    \def\secondColumn{#2}
    \setcolumnwidth{\fill, 2.4 cm}
    \begin{paracol}{2}
}{
    \switchcolumn \raggedleft \secondColumn
    \end{paracol}
    \endonecolentry
} % new environment for two column entries

\newenvironment{header}{
    \setlength{\topsep}{0pt}\par\kern\topsep\centering\linespread{1.5}
}{
    \par\kern\topsep
} % new environment for the header

\newcommand{\placelastupdatedtext}{% \placetextbox{<horizontal pos>}{<vertical pos>}{<stuff>}
  \AddToShipoutPictureFG*{% Add <stuff> to current page foreground
    \put(
        \LenToUnit{\paperwidth-2 cm-0.2 cm+0.05cm},
        \LenToUnit{\paperheight-1.0 cm}
    ){\vtop{{\null}\makebox[0pt][c]{
        \small\color{gray}\textit{}\hspace{\widthof{Last updated in September 2024}}
    }}}%
  }%
}%

% save the original href command in a new command:
\let\hrefWithoutArrow\href

% new command for external links:
\renewcommand{\href}[2]{\hrefWithoutArrow{#1}{\ifthenelse{\equal{#2}{}}{ }{#2 }\raisebox{.15ex}{\footnotesize \faExternalLink*}}}


\begin{document}
    \newcommand{\AND}{\unskip
        \cleaders\copy\ANDbox\hskip\wd\ANDbox
        \ignorespaces
    }
    \newsavebox\ANDbox
    \sbox\ANDbox{}

    \placelastupdatedtext
    \begin{header}
        \textbf{\fontsize{24 pt}{24 pt}\selectfont Hao-Hsuan Chang}

        \vspace{0.2 cm}

        \normalsize
        
        \mbox{\hrefWithoutArrow{mailto:william29181@gmail.com}{\color{black}{\footnotesize\faEnvelope[regular]}\hspace*{0.13cm}haohsuanchang2918@gmail.com}}%ha
        \kern 0.25 cm%
        \AND%
        \kern 0.25 cm%
        \mbox{\hrefWithoutArrow{https://haohsuan2918.github.io/}{\color{black}{\footnotesize\faLink}\hspace*{0.13cm}Homepage}}%
        \kern 0.25 cm%
        \AND%
        \kern 0.25 cm%
        \mbox{\hrefWithoutArrow{https://scholar.google.com/citations?user=Ra2fZSAAAAAJ&hl=en}{\color{black}{\faGraduationCap}\hspace*{0.13cm}Google Scholar}}%
        \kern 0.25 cm%
        \AND%
        \kern 0.25 cm%
        \mbox{\hrefWithoutArrow{https://www.linkedin.com/in/hao-hsuan-chang-8862b7170}{\color{black}{\footnotesize\faLinkedinIn}\hspace*{0.13cm}Linkedin}}%
        \kern 0.25 cm%
        \AND%
        \kern 0.25 cm%
        \mbox{\hrefWithoutArrow{https://github.com/haohsuan2918}{\color{black}{\footnotesize\faGithub}\hspace*{0.13cm}Github}}%
        
        


        

        
    \end{header}

    \vspace{0.2 cm - 0.2 cm}


    \section{Professional Summary}
    I am currently a Staff Research Engineer at Samsung Research America, specializing in AI solutions for wireless sensing and communication systems. With over 8 years of experience, I have developed a broad range of low-complexity AI applications optimized for commercial embedded devices, with expertise in time-series prediction, reinforcement learning, recurrent neural networks, and LLM finetuning. I have a strong track record of driving projects from concept to deployment, resulting in high-impact publications, patents, and commercial products.

    
    \section{Experience}

        \begin{twocolentry}{  
            
        {\small \textit{Mar 2025–Present}}
        
        {\small \textit{Sep 2021–Feb 2025}}}

            \textbf{Staff Research Engineer}, \textit{Samsung Research America}

            \textbf{Senior Research Engineer}, \textit{Samsung Research America}

        \end{twocolentry}

        \vspace{0.10 cm}
        \begin{onecolentry}
            \begin{highlights}
                \item Designed AI models for time-series–based sleep stage prediction using wireless signals, leveraging residual learning, transformers, and contrastive learning techniques; achieved accuracy comparable to clinical sleep study, with commercialization planned for 2027 (2 publications, 4 patents).

                \item Fine-tuned LLMs using predicted sleep stages and contextual data to deliver personalized sleep coaching and actionable health insights, enabling human-like interaction and improving user engagement.

                
                \item Optimized and deployed AI models on embedded devices by applying CNN quantization and TensorFlow Lite optimization for low-latency inference, and converting Python-based models to C\texttt{++} with multi-threading and optimized data pipelines to achieve real-time sleep stage prediction on commercial hardware.

                \item Designed an AI model to detect human count in different areas using CBRS wireless signals; presented a demo at the FIFA World Cup fan festival at Dallas Fair Park; built a multi-camera monitoring system using homographies to collect training data ground truth.
                
                % Reduced model complexity through CNN quantization techniques and TensorFlow Lite optimization, enabling compact AI models suitable for deployment on commercial embedded devices.


                % \item Converted Python-based AI models to C\texttt{++} for embedded deployment, integrating multi-threading and optimizing the data preprocessing pipeline to enable real-time sleep stage inference on embedded devices.
            \end{highlights}
        \end{onecolentry}


        \vspace{0.10 cm}

        \begin{twocolentry}{

        {\small \textit{Aug 2019–Dec 2019}}}
            \textbf{Research Intern}, \textit{Samsung Research America}

        \end{twocolentry}

        \vspace{0.10 cm}
        \begin{onecolentry}
            \begin{highlights}
                \item Designed deep reinforcement learning-based load balancing optimization (1 publication).
            \end{highlights}
        \end{onecolentry}

        \vspace{0.10 cm}

        \begin{twocolentry}{

        {\small \textit{Aug 2017–Aug 2021}}}
            \textbf{Graduate Research Assistant}, \textit{Virginia Tech}

        \end{twocolentry}

        \vspace{0.10 cm}
        \begin{onecolentry}
            \begin{highlights}
                \item Designed distributive multi-agent deep reinforcement learning and lightweight recurrent neural network architectures for time-series prediction of multi-user spectrum access behavior (7 publications).
            \end{highlights}
        \end{onecolentry}



    
    \section{Education}
        \begin{twocolentry}{
        {\small \textit{Aug 2017–Aug 2021}}}
            \makebox[1.5cm]{Ph.D}   \textbf{Virginia Tech}, Electrical and Computer Engineering 
        \end{twocolentry}

        \vspace{0.05 cm}

        \begin{twocolentry}{
        {\small \textit{Sep 2013–Jun 2015}}}
            \makebox[1.5cm]{M.S}    \textbf{National Taiwan University}, Communication Engineering
        \end{twocolentry}

        \vspace{0.05 cm}

        \begin{twocolentry}{
        {\small \textit{Sep 2008–Jun 2013}}}
            \makebox[1.5cm]{B.S}   \textbf{National Taiwan University}, Electrical Engineering 
        \end{twocolentry}

        % \vspace{0.10 cm}

        
    \section{Technical Skills}

        \begin{onecolentry}
            \textbf{Certification:} AWS Certified Machine Learning Engineer, AWS Certified Cloud Practitioner
        \end{onecolentry}

        \vspace{0.1 cm}
        
        \begin{onecolentry}
            \textbf{Programming:} Python, C\texttt{++}, SQL, Matlab
        \end{onecolentry}

        \vspace{0.1 cm}

        \begin{onecolentry}
            \textbf{MLOps \& Tools:} PyTorch, TensorFlow, ONNX, Kubernetes
        \end{onecolentry}

        \vspace{0.1 cm}

        \begin{onecolentry}
            \textbf{Deep Learning:} LLM finetuning, Transformers, Reinforcement Learning, RNNs, CNNs, XGBoost
        \end{onecolentry}







    

    \section{Selected Publications}



        
        \begin{samepage}
        
            \begin{twocolentryPub}{
                {\footnotesize IoT 2018}
            }
                {\footnotesize \textbf{Distributive Dynamic Spectrum Access through Deep Reinforcement Learning} } [270+ citations]
            \end{twocolentryPub}

            \vspace{0.05 cm}

            \begin{twocolentryPub}{
                {\footnotesize WCL 2019}
            }
                {\footnotesize \textbf{Deep Residual Learning Meets OFDM Channel Estimation}} [200+ citations]
            \end{twocolentryPub}

            \vspace{0.05 cm}

            \begin{twocolentryPub}{
                {\footnotesize TNNLS 2020}
            }
                {\footnotesize \textbf{DEQN and Its Application in Dynamic Spectrum Sharing for 5G and Beyond}} [80+ citations]
            \end{twocolentryPub}

            \vspace{0.05 cm}

            \begin{twocolentryPub}{
                {\footnotesize TWC 2020}
            }
                {\footnotesize \textbf{Learning for detection: MIMO-OFDM symbol detection through downlink pilots}} [70+ citations]
            \end{twocolentryPub}

            \vspace{0.05 cm}

            \begin{twocolentryPub}{
                {\footnotesize ICASSP 2024}
            }
                {\footnotesize \textbf{Multi-Person Respiration Rate Estimation with Single Pair of Transmit and Receive Antenna}}
            \end{twocolentryPub}

            \vspace{0.05 cm}

            % \begin{twocolentryPub}{
            %     {\footnotesize TWC 2024}
            % }
            %     {\footnotesize \textbf{Efficient Deep Reinforcement Learning for Partially Observable Dynamic Spectrum Access}}
            % \end{twocolentryPub}

            % \vspace{0.05 cm}

            % \begin{twocolentryPub}{
            %     {\footnotesize TWC 2024}
            % }
            %     {\footnotesize \textbf{Deep Reinforcement Learning for Dynamic Spectrum Access: Convergence Analysis}}
            % \end{twocolentryPub}

            % \vspace{0.05 cm}


            \begin{twocolentryPub}{
                {\footnotesize TWC 2023}
            }
                {\footnotesize \textbf{Federated Multi-Agent Deep Reinforcement Learning for Dynamic Spectrum Access}}
            \end{twocolentryPub}

            \vspace{0.05 cm}


            \begin{twocolentryPub}{
                {\footnotesize ToN 2022}
            }
                {\footnotesize \textbf{Decentralized Deep Reinforcement Learning Meets Mobility Load Balancing}}
            \end{twocolentryPub}

            \vspace{0.05 cm}

            % \begin{twocolentryPub}{
            %     {\footnotesize MILCOM 2022}
            % }
            %     {\footnotesize \textbf{MADRL Based Scheduling for 5G and Beyond}}
            % \end{twocolentryPub}

            % \vspace{0.05 cm}






        \end{samepage}



\end{document}